\clearpage
\section{Prototipos de alta fidelidad}

Los cinco prototipos recorren escenarios complementarios de Karisma Data: consulta,
análisis, toma de decisiones, publicación y gobierno. Cada uno responde a una tarea y a una
necesidad de información distinta, mientras el sistema visual permanece estable para reducir
el reaprendizaje y reforzar la identidad del producto.

\subsection{Análisis y exportación de liquidez}

El recorrido de liquidez parte de la localización de información, continúa con la revisión de
indicadores y filtros, y culmina en la preparación de una exportación. Esta secuencia permite
pasar de una vista general a evidencia utilizable sin abandonar el contexto del análisis.

Sus ocho pantallas hacen visible el avance desde el inicio personalizado hasta la confirmación
del archivo generado. En cada etapa se conserva el origen de los datos, las variables elegidas
y las restricciones de acceso, de modo que la persona pueda revisar sus decisiones antes de
exportar.

\figuraaiv[0.94\linewidth]{figuras/a4/secuencia_flujo_1.png}{Secuencia completa del prototipo de análisis y exportación de liquidez: inicio, exploración, detalle, selección, filtros, vista previa, configuración y confirmación.}{fig:a4-secuencia1}

\figuraaiv{figuras/a4/flujo_1_variables.png}{Selección de variables dentro del flujo de análisis y exportación de liquidez. La pantalla muestra el paso activo, las restricciones de acceso y el resumen de la selección.}{fig:a4-flujo1}

\subsection{Señal de riesgo para la Junta}

A diferencia del recorrido anterior, este escenario se concentra en una señal de riesgo y la
acompaña desde su detección hasta la consulta de su contexto. La jerarquía privilegia el nivel
de riesgo, la procedencia del dato y la evidencia que respalda una decisión ejecutiva.

Las cinco pantallas conducen de la lectura panorámica a un resumen listo para presentar ante
la Junta. El recorrido no se limita a mostrar una alerta: primero solicita contexto, después
expone la explicación y la composición del indicador, y finalmente sintetiza los hallazgos.

\figuraaiv[0.86\linewidth]{figuras/a4/secuencia_flujo_2.png}{Secuencia completa del prototipo de señal de riesgo: panorama, detección, explicación, composición y resumen ejecutivo.}{fig:a4-secuencia2}

\figuraaiv[0.64\linewidth]{figuras/a4/flujo_2_riesgo.png}{Cambio detectado dentro del flujo de señal de riesgo. El estado combina magnitud, umbral, evolución y acción siguiente.}{fig:a4-flujo2}

\subsection{Publicación de una versión del catálogo}

La misma lógica de trazabilidad se traslada a la revisión y publicación controlada de una
versión del catálogo. Los pasos visibles, las confirmaciones y la retroalimentación del
sistema reducen la posibilidad de publicar cambios sin comprender su alcance.

El recorrido distribuye la tarea en seis momentos: selección del elemento, documentación del
cambio, comparación entre versiones, análisis de impacto, revisión final y confirmación. Esta
separación permite conservar evidencia y anticipar dependencias antes de publicar.

\figuraaiv[0.94\linewidth]{figuras/a4/secuencia_flujo_3.png}{Secuencia completa del prototipo de publicación del catálogo, desde la selección del registro hasta la confirmación de la nueva versión.}{fig:a4-secuencia3}

\figuraaiv{figuras/a4/flujo_3_versiones.png}{Comparación de versiones del catálogo. La interfaz contrasta definición, regla histórica, estructura e interpretación antes de continuar.}{fig:a4-flujo3}

\subsection{Administración de accesos}

En materia de gobierno, el flujo de accesos permite revisar y modificar permisos según la
función de cada persona. La interfaz expone el alcance de cada autorización y reserva las
confirmaciones para acciones que afectan a otras personas usuarias.

Las cinco pantallas aplican el principio de permiso mínimo suficiente: parten de una solicitud,
revisan su justificación, acotan las capacidades habilitadas y conservan un registro de la
autorización. La confirmación final comunica vigencia, alcance y evidencia disponible.

\figuraaiv[0.86\linewidth]{figuras/a4/secuencia_flujo_4.png}{Secuencia completa del prototipo de administración de accesos: solicitud, revisión, definición del permiso mínimo, autorización y alta.}{fig:a4-secuencia4}

\figuraaiv[0.64\linewidth]{figuras/a4/flujo_4_autorizacion.png}{Confirmación de autorización. La pantalla conserva la justificación, responsables, vigencia, alcance, exclusiones y controles verificados.}{fig:a4-flujo4}

\subsection{Herramienta integral}

Por último, el centro de trabajo articula las capacidades principales del portal y hace
visible la continuidad entre navegación, consulta, análisis y gobierno. El recorrido presenta
Karisma Data como un sistema integrado, en lugar de una suma de módulos independientes.

Sus seis vistas comprueban que la navegación común puede sostener tareas de exploración,
exportación, monitoreo, administración e integración sin perder el contexto. Más que una tarea
lineal, este prototipo representa la arquitectura global del producto y sus principales puntos
de entrada.

\figuraaiv[0.94\linewidth]{figuras/a4/secuencia_flujo_5.png}{Secuencia completa del prototipo integral: centro de trabajo, análisis, exportación, riesgos, gobierno e integración.}{fig:a4-secuencia5}

\figuraaiv{figuras/a4/flujo_5_centro_trabajo.png}{Centro de trabajo del flujo integral. La vista reúne indicadores, módulos principales y actividad reciente bajo una navegación común.}{fig:a4-flujo5}

\clearpage
